\chapter{The Linear Programming Category}\label{sec:glp}

The theory developed so far has extended the categorical framework
by adding a new class of  morphisms. We now introduce the
final and most expressive theory of the thesis:
a  graphical algebra for parametric linear programs.
Linear programs combine  both the inequality structure of polyhedra and the notion of an objective function,
and their value functions are precisely the polyhedral convex  functions.
The central result of this section is the presentation of a
complete monoidal theory for right-hand side parametric linear programming.

\section{Semantics: the prop \texorpdfstring{$\mathbf{LPRel}$}{LPRel}}

We begin with several classical results that establish the tight
relationship between piecewise affine convex functions, polyhedral
functions, and parametric linear programs. These will serve as the
semantic foundation for the graphical algebra introduced in the following
section, and they justify describing the morphisms of this chapter's
semantic category, $\mathbf{LPRel}$, interchangeably as parametric linear
programs or as polyhedral functions.

\begin{definition}[Piecewise Affine Convex Function]
A function $f : \mathbb{R}^n \to \mathbb{R}$ is called
\emph{piecewise affine convex} if there exist finitely many affine
functions $a_i^\top x + b_i$, $i = 1,\dots,k$, such that
\[
f(x) = \max_{i=1,\dots,k} \{ a_i^\top x + b_i \}.
\]
\end{definition}

\begin{definition}[Polyhedral Function]
A function $f : \mathbb{R}^n \to \mathbb{R} \cup \{+\infty\}$ is
\emph{polyhedral} if its epigraph
\[
\mathrm{epi}(f) := \{ (x,t) \mid t \ge f(x) \}
\]
is a polyhedron.
\end{definition}

\begin{definition}[RHS Parametric Linear Program]
A \emph{right-hand side parametric linear program} has value function
\[
v(b) := \inf_{x} \{ c^\top x \mid Ax \ge b \}.
\]
\end{definition}

The following two propositions establish that all three notions above
are equivalent descriptions of the same class of functions. This
equivalence is the semantic cornerstone of the section: it is what
allows us to treat polyhedral functions and parametric linear programs
as morphisms in the same category.

\begin{proposition}
A function is piecewise affine convex if and only if it is a polyhedral
function.
\end{proposition}

\begin{proof}
\textbf{($\Rightarrow$)} Let
$f(x) = \max_{i=1,\dots,k} (a_i^\top x + b_i)$.
Then
\[
\mathrm{epi}(f) = \bigcap_{i=1}^k \{ (x,t) \mid t \ge a_i^\top x + b_i \},
\]
which is an intersection of finitely many halfspaces, hence a polyhedron.

\medskip
\textbf{($\Leftarrow$)} If $f$ is polyhedral, its epigraph is a
polyhedron, so there exist finitely many affine inequalities
$t \ge a_i^\top x + b_i$ such that
$f(x) = \max_{i=1,\dots,k} (a_i^\top x + b_i)$.
Thus $f$ is piecewise affine and convex.
\end{proof}

\begin{proposition}
A function $v : \mathbb{R}^m \to \mathbb{R} \cup \{+\infty\}$ is
piecewise affine convex if and only if it is the value function of a
right-hand side parametric linear program.
\end{proposition}

\begin{proof}
\textbf{($\Rightarrow$)} Let
$v(b) = \max_{i=1,\dots,k} (y^i)^\top b$
and define $Y := \mathrm{conv}\{y^1,\dots,y^k\}$.
Then $v(b) = \max_{y \in Y} b^\top y$.
Since $Y$ is a polyhedron admitting the representation
$Y = \{ y \mid A^\top y = c,\; y \ge 0 \}$,
linear programming duality gives
$v(b) = \inf \{ c^\top x \mid Ax \ge b \}$.

\medskip
\textbf{($\Leftarrow$)} Let
$v(b) = \inf \{ c^\top x \mid Ax \ge b \}$.
Dualising,
$v(b) = \sup \{ b^\top y \mid A^\top y = c,\; y \ge 0 \}$.
Since the feasible set $Y = \{ y \mid A^\top y = c,\; y \ge 0 \}$
is a polyhedron, the supremum is attained at its extreme points
$y^1,\dots,y^k$, yielding
$v(b) = \max_{i=1,\dots,k} (y^i)^\top b$.
Thus $v$ is piecewise affine convex.
\end{proof}

The following proposition is the key technical lemma for the completeness
proof of $\mathbf{GLP}$. It states that two parametric linear programs
with the same value function must share the same epigraph polyhedron,
providing the bridge between semantic equality and polyhedral equality
that the completeness argument requires.

\begin{proposition}\label{prop:equalepi}
Let
\[
v_i(y) := \inf\{\mathbf{1}^\top x \mid A_i x \ge B_i y\},
\qquad i=1,2,
\]
where $A_i \in \mathbb{R}^{m_i \times n_i}$ and
$B_i \in \mathbb{R}^{m_i \times k}$. If $v_1(y) = v_2(y)$ for all
$y \in \mathbb{R}^k$, then both programs admit a common representation
through the same epigraph polyhedron:
\[
v_i(y) = \inf\{t \mid (y,t) \in E\},
\qquad
E := \mathrm{epi}(v_1) = \mathrm{epi}(v_2).
\]
\end{proposition}

\begin{proof}
For each $i = 1,2$, the epigraph of $v_i$ satisfies
\[
(y,t) \in \mathrm{epi}(v_i)
\iff
\exists x : A_i x \ge B_i y,\ \mathbf{1}^\top x \le t,
\]
so $\mathrm{epi}(v_i)$ is the projection onto $(y,t)$ of the polyhedron
$P_i = \{(x,y,t) \mid A_i x \ge B_i y,\ \mathbf{1}^\top x \le t\}$.
Projections of polyhedra are polyhedra, so each $\mathrm{epi}(v_i)$ is
a polyhedron. Since $v_1 = v_2$ pointwise, we have
$\mathrm{epi}(v_1) = \mathrm{epi}(v_2) =: E$.
Setting $E = \{(y,t) \mid Cy + dt \ge 0\}$, we obtain
$v_i(y) = \inf\{t \mid Cy + dt \ge 0\}$ for both $i = 1, 2$.
\end{proof}

With the equivalences above in place, we now define the semantic
category of this chapter. $\mathbf{LPRel}$ describes morphisms as
parametric linear programs; by the propositions above, these coincide
with polyhedral functions, and we will use both descriptions
interchangeably according to which is more convenient.

\begin{definition}[The prop $\mathbf{LPRel}$ of Linear Programs]
The prop $\mathbf{LPRel}$ is defined as follows:
\begin{itemize}
    \item Objects are natural numbers $n \in \mathbb{N}$, interpreted
    as the space $\mathbb{R}^n$.
    \item A morphism $n \to m$ is a right-hand side parametric linear
    program with value function
    \[
    (x,y) \mapsto \inf_z \{ c^\top z \mid Az \ge B[x\ y]^\top \},
    \]
    where $x \in \mathbb{R}^n$ and $y \in \mathbb{R}^m$.
    \item Composition is given by minimisation over the shared variable.
    \item The monoidal product on objects is addition $n \otimes m := n+m$,
    and on morphisms is direct sum.
\end{itemize}
\end{definition}

\section{Syntax: the algebra \texorpdfstring{$\mathbf{GLP}$}{GLP}}

With the semantic equivalences established, we now introduce the final
graphical algebra of the thesis. The algebra $\mathbf{GLP}$ extends
$\mathbf{GPA}$ with a single new generator whose semantics encodes the
linear objective function $c^\top x$. Just as the gray arrow in
$\mathbf{GQA}$ lifted the theory from binary relations to weighted
relations by "measuring" the $\ell_2$ norm, this new generator lifts $\mathbf{GPA}$
from polyhedral constraints to full linear programs by attaching a
linear cost to the feasible region.

\begin{definition}[$\mathbf{GLP}$]
    We call $\mathbf{GLP}$ the symmetric monoidal theory obtained by
    extending $\mathbf{GPA}$ with the following generator:
    \begin{figure}[H]
        \centering
        \begin{tikzpicture}
    \node at (0.2,0.5) {$\llbracket$};
    \pic (g1) at (0.500,0.500) {linear};
    \node at (1.25,0.5) {$\rrbracket$};
    \node at (1.65,0.5) {$=$};
    \node[anchor=west] at (1.9,0.5) {$(\bullet,x) \mapsto x$};
\end{tikzpicture}

    \end{figure}
    together with the following axioms governing its interaction with
    the existing generators:
    \begin{figure}[H]
        \centering
        \begin{tabular}{r c c c}
  \diagrow{(linearity)}{
      \pic (g1) at (0.500,1.000) {linear};
      \pic (g2) at (0.500,0.000) {linear};
  }{
      \pic (g1) at (0.500,0.500) {linear};
      \pic (g2) at (1.750,0.500) {comultiplication};
      \draw (g1-out-0) to[out=0,in=180] (g2-in-0);
  }
  \diagrow{(zeroeq)}{
      \pic (g1) at (0.500,0.500) {linear};
      \begin{scope}[shift={(1.250,0.000)}]
      \pic (g2) at (0.500,0.500) {cozero};
      \end{scope}
      \draw (g1-out-0) to[out=0,in=180] (g2-in-0);
  }{
      \node[draw, dotted] {};
  }
  \diagrow{(inf)}{
      \pic (g1) at (0.500,0.500) {linear};
      \begin{scope}[shift={(1.250,0.000)}]
      \pic (g2) at (0.500,0.500) {ge};
      \end{scope}
      \draw (g1-out-0) to[out=0,in=180] (g2-in-0);
  }{
      \pic (g1) at (0.500,0.500) {linear};
  }
\end{tabular}
    \end{figure}
\end{definition}

The axioms represent expected properties of the linear cost function.
The one on the left express the linear property $f(x) + f(y) = f(x+y)$.
Observe that $f$ is the interpretation of the $\mushroom$ and the
sum denotes the vertical composition; composing two $\mushroom$ vertically
(adding) is the same as adding the wires in a single $\mushroom$. This
axiom is also the fundamental property of LPs, which is the fact that
they are invariant to column stochastic matrices
(see~\ref{rmk:stochastic}). In turn, the
equation on the right express the fundamental property of
minimization, the infimum of any quantity lies in its lower bound.
\[\inf_x \{ x | x \ge y \} =  y\]

Finally, the axiom at the middle is
the tautological case. \[inf 0 = 0\]

\begin{remark}[Invariance under column-stochastic matrices]\label{rmk:stochastic}
    The \textit{linearity} axiom encodes the fundamental invariance of linear
    programs under column-stochastic reweightings of the objective: for any
    $n \times m$ matrix $S$ whose columns each sum to one,
    \[
        \inlinegen{linear}^{\otimes n} ; \; S^{op} \;=\; \inlinegen{linear}^{\otimes m},
    \]
    where $S^{op}$ denotes the comatrix of $S$. Note that only the column sums are
    used; non-negativity of the entries plays no role. The proof is diagrammatic,
    we display the proof for a
    $2 \times 2$ matrix $S$ with $s_{1,1} + s_{2,1} = s_{1,2} + s_{2,2} = 1$;
    every step is an instance of an axiom scheme that applies wire-by-wire, so
    the same chain of rewrites goes through verbatim for arbitrary $n \times m$.
    \begin{figure}[H]
        \centering
        \begin{center}
    \begin{tikzpicture}[axpic]
      \pic (l1) at (0.500,0.500) {linear};
      \pic[val=S] (m1) at (1.750,0.500) {comatrix};
      \draw (l1-out-0) to[out=0,in=180] (m1-in-0);
    \end{tikzpicture}
    \;\begin{tikzpicture}[axpic]\node{$\overset{(def)}{=}$};\end{tikzpicture}\;
    \begin{tikzpicture}[axpic]
      \pic (l1) at (0.500,2.500) {linear};
      \pic (l2) at (0.500,0.500) {linear};
      \pic (c1) at (1.750,2.500) {comultiplication};
      \pic (c2) at (1.750,0.500) {comultiplication};
      \pic[val=s_{1,1}] (s11) at (3.000,3.000) {coscalar};
      \pic[val=s_{1,2}] (s12) at (3.000,2.000) {coscalar};
      \pic[val=s_{2,1}] (s21) at (3.000,1.000) {coscalar};
      \pic[val=s_{2,2}] (s22) at (3.000,0.000) {coscalar};
      \pic (k1) at (5.000,2.500) {cocopy};
      \pic (k2) at (5.000,0.500) {cocopy};
      \draw (l1-out-0) to[out=0,in=180] (c1-in-0);
      \draw (l2-out-0) to[out=0,in=180] (c2-in-0);
      \draw (c1-out-0) to[out=0,in=180] (s11-in-0);
      \draw (c1-out-1) to[out=0,in=180] (s12-in-0);
      \draw (c2-out-0) to[out=0,in=180] (s21-in-0);
      \draw (c2-out-1) to[out=0,in=180] (s22-in-0);
      \draw (s11-out-0) to[out=0,in=180] (k1-in-0);
      \draw (s21-out-0) to[out=0,in=180] (k1-in-1);
      \draw (s12-out-0) to[out=0,in=180] (k2-in-0);
      \draw (s22-out-0) to[out=0,in=180] (k2-in-1);
    \end{tikzpicture} \\
    \vspace{2mm}

    \;\begin{tikzpicture}[axpic]\node{$\overset{(linearity)}{=}$};\end{tikzpicture}\;
    \begin{tikzpicture}[axpic]
      \pic (l1) at (0.500,1.500) {linear};
      \pic (c0) at (1.750,1.500) {comultiplication};
      \pic (c1) at (3.000,2.500) {comultiplication};
      \pic (c2) at (3.000,0.500) {comultiplication};
      \pic[val=s_{1,1}] (s11) at (4.250,3.000) {coscalar};
      \pic[val=s_{1,2}] (s12) at (4.250,2.000) {coscalar};
      \pic[val=s_{2,1}] (s21) at (4.250,1.000) {coscalar};
      \pic[val=s_{2,2}] (s22) at (4.250,0.000) {coscalar};
      \pic (k1) at (6.250,2.500) {cocopy};
      \pic (k2) at (6.250,0.500) {cocopy};
      \draw (l1-out-0) to[out=0,in=180] (c0-in-0);
      \draw (c0-out-0) to[out=0,in=180] (c1-in-0);
      \draw (c0-out-1) to[out=0,in=180] (c2-in-0);
      \draw (c1-out-0) to[out=0,in=180] (s11-in-0);
      \draw (c1-out-1) to[out=0,in=180] (s12-in-0);
      \draw (c2-out-0) to[out=0,in=180] (s21-in-0);
      \draw (c2-out-1) to[out=0,in=180] (s22-in-0);
      \draw (s11-out-0) to[out=0,in=180] (k1-in-0);
      \draw (s21-out-0) to[out=0,in=180] (k1-in-1);
      \draw (s12-out-0) to[out=0,in=180] (k2-in-0);
      \draw (s22-out-0) to[out=0,in=180] (k2-in-1);
    \end{tikzpicture} \\
    \vspace{2mm}

    \;\begin{tikzpicture}[axpic]\node{$\overset{(coassoc,\,cocomm)}{=}$};\end{tikzpicture}\;
    \begin{tikzpicture}[axpic]
      \pic (l1) at (0.500,1.500) {linear};
      \pic (c0) at (1.750,1.500) {comultiplication};
      \pic (c1) at (3.000,2.500) {comultiplication};
      \pic (c2) at (3.000,0.500) {comultiplication};
      \pic[val=s_{1,1}] (s11) at (4.250,3.000) {coscalar};
      \pic[val=s_{2,1}] (s21) at (4.250,2.000) {coscalar};
      \pic[val=s_{1,2}] (s12) at (4.250,1.000) {coscalar};
      \pic[val=s_{2,2}] (s22) at (4.250,0.000) {coscalar};
      \pic (k1) at (5.500,2.500) {cocopy};
      \pic (k2) at (5.500,0.500) {cocopy};
      \draw (l1-out-0) to[out=0,in=180] (c0-in-0);
      \draw (c0-out-0) to[out=0,in=180] (c1-in-0);
      \draw (c0-out-1) to[out=0,in=180] (c2-in-0);
      \draw (c1-out-0) to[out=0,in=180] (s11-in-0);
      \draw (c1-out-1) to[out=0,in=180] (s21-in-0);
      \draw (c2-out-0) to[out=0,in=180] (s12-in-0);
      \draw (c2-out-1) to[out=0,in=180] (s22-in-0);
      \draw (s11-out-0) to[out=0,in=180] (k1-in-0);
      \draw (s21-out-0) to[out=0,in=180] (k1-in-1);
      \draw (s12-out-0) to[out=0,in=180] (k2-in-0);
      \draw (s22-out-0) to[out=0,in=180] (k2-in-1);
    \end{tikzpicture} \\
    \vspace{2mm}

    \;\begin{tikzpicture}[axpic]\node{$\overset{(s.add)}{=}$};\end{tikzpicture}\;
    \begin{tikzpicture}[axpic]
      \pic (l1) at (0.500,1.000) {linear};
      \pic (c0) at (1.750,1.000) {comultiplication};
      \node[boxnode_left] (s1) at (3.750,1.500) {$s_{1,1}{+}s_{2,1}$};
      \node[boxnode_left] (s2) at (3.750,0.500) {$s_{1,2}{+}s_{2,2}$};
      \draw (l1-out-0) to[out=0,in=180] (c0-in-0);
      \draw (c0-out-0) to[out=0,in=180] (s1.west);
      \draw (c0-out-1) to[out=0,in=180] (s2.west);
      \draw (s1.east) -- +(0.400,0);
      \draw (s2.east) -- +(0.400,0);
    \end{tikzpicture}
    \;\begin{tikzpicture}[axpic]\node{$\overset{(stoch.)}{=}$};\end{tikzpicture}\;
    \begin{tikzpicture}[axpic]
      \pic (l1) at (0.500,1.000) {linear};
      \pic (c0) at (1.750,1.000) {comultiplication};
      \draw (l1-out-0) to[out=0,in=180] (c0-in-0);
    \end{tikzpicture}
    \;\begin{tikzpicture}[axpic]\node{$\overset{(linearity)}{=}$};\end{tikzpicture}\;
    \begin{tikzpicture}[axpic]
      \pic (l1) at (0.500,1.500) {linear};
      \pic (l2) at (0.500,0.500) {linear};
    \end{tikzpicture}
\end{center}

    \end{figure}
\end{remark}

\begin{definition}
    We call $\mathrm{FreeP}_{\mathbf{GLP}}$ the PROP freely generated
    by $\mathbf{GLP}$.
\end{definition}

The normal form theorem for $\mathbf{GLP}$ extends the one for
$\mathbf{GPA}$: every diagram can be reduced to a $\mathbf{GPA}$ block
encoding the constraints, composed with the new generator encoding the
objective. The proof proceeds by structural induction, using the existing
normal form for $\mathbf{GPA}$ as the inductive base.

\begin{theorem}[Normal Form]\label{thm:NormalFormGLP}
    For every diagram $c \in \mathbf{GLP}$, there exists a diagram $c'$
    in normal form such that $c = c'$:
    \begin{figure}[H]
        \centering
        \begin{tikzpicture}
    \pic (g1) at (0.500,1.300) {linear};
    \coordinate (id_in_2) at (0,0.700);
    \coordinate (id_out_3) at (1.000,0.700);
    \draw (id_in_2) -- (id_out_3);
    \pic[val=A] (g4) at (2.000,1.000) {2drelation};
    \draw (g1-out-0) to[out=0,in=180] (g4-in-0);
    \draw (id_out_3) to[out=0,in=180] (g4-in-1);
    \pic (g5) at (3.500,1.000) {ge};
    \draw (g4-out-0) to[out=0,in=180] (g5-in-0);
    \pic[val=B] (g6) at (5.000,1.000) {corelation};
    \draw (g5-out-0) to[out=0,in=180] (g6-in-0);
\end{tikzpicture}

    \end{figure}
    With $A$ and $b$ being an affine transformation. Namely:
    \begin{figure}[H]
        \centering
        \begin{center}
    \begin{tikzpicture}[axpic]
      \pic[val=A'] (g1) at (1.000,1.500) {matrix};
      \pic (g2) at (1.500,0.500) {affine};
      \pic (g3) at (3.000,1.000) {multiplication};
      \draw (g1-out-0) to[out=0,in=180] (g3-in-0);
      \draw (g2-out-0) to[out=0,in=180] (g3-in-1);
    \end{tikzpicture}
    \;\begin{tikzpicture}[axpic]\node{$=$};\end{tikzpicture}\;
    \begin{tikzpicture}[axpic]
      \pic[val=A] (g1) at (1.000,0.500) {relation};
    \end{tikzpicture}
\end{center}

    \end{figure}
\end{theorem}

\begin{proof}
    The proof is by structural induction. First we show the two base cases:
    \begin{itemize}
        \item The generator $\mushroom$
        \begin{figure}[H]
            \centering
            \begin{center}
    \begin{tikzpicture}[axpic]
      \pic (g1) at (0.500,0.500) {linear};
    \end{tikzpicture}
    \;\begin{tikzpicture}[axpic]\node{$\overset{(inf)}{=}$};\end{tikzpicture}\;
    \begin{tikzpicture}[axpic]
      \pic (g1) at (0.500,0.500) {linear};
      \pic (g2) at (1.750,0.500) {ge};
      \draw (g1-out-0) to[out=0,in=180] (g2-in-0);
    \end{tikzpicture} \\

    \;\begin{tikzpicture}[axpic]\node{$\overset{(identity)}{=}$};\end{tikzpicture}\;
    \begin{tikzpicture}[axpic]
      \pic (g1) at (0.500,0.500) {linear};
      \pic[val=I] (g2) at (2.000,0.500) {matrix};
      \pic (g3) at (3.500,0.500) {ge};
      \pic[val=I] (g4) at (5.000,0.500) {comatrix};
      \draw (g1-out-0) to[out=0,in=180] (g2-in-0);
      \draw (g2-out-0) to[out=0,in=180] (g3-in-0);
      \draw (g3-out-0) to[out=0,in=180] (g4-in-0);
    \end{tikzpicture} \\
    \vspace{2mm}

    \;\begin{tikzpicture}[axpic]\node{$\overset{(unity)}{=}$};\end{tikzpicture}\;
    \begin{tikzpicture}[axpic]
      \pic (g1) at (0.500,1.000) {linear};
      \pic[val=I] (g2) at (2.000,1.000) {matrix};
      \pic (z1) at (2.500,0.000) {zero};
      \pic (m3) at (3.500,0.500) {multiplication};
      \pic (g3) at (4.500,0.500) {ge};
      \pic[val=I] (g4) at (6.000,0.500) {comatrix};
      \draw (g1-out-0) to[out=0,in=180] (g2-in-0);
      \draw (g2-out-0) to[out=0,in=180] (m3-in-0);
      \draw (z1-out-0) to[out=0,in=180] (m3-in-1);
      \draw (m3-out-0) to[out=0,in=180] (g3-in-0);
      \draw (g3-out-0) to[out=0,in=180] (g4-in-0);
    \end{tikzpicture} \\
    \vspace{2mm}

    \;\begin{tikzpicture}[axpic]\node{$=$};\end{tikzpicture}\;
    \begin{tikzpicture}[axpic]
      \pic (g1) at (0.500,0.300) {linear};
      \coordinate (id_in_2) at (0.500,-0.300);
      \coordinate (id_out_3) at (1.000,-0.300);
      \draw (id_in_2) -- (id_out_3);
      \node[left] at (id_in_2) {$0$};
      \pic[val=I] (g2) at (2.000,0.000) {2drelation};
      \pic (g3) at (3.500,0.000) {ge};
      \pic[val=I] (g4) at (5.000,0.000) {comatrix};
      \draw (g1-out-0) to[out=0,in=180] (g2-in-0);
      \draw (id_out_3) to[out=0,in=180] (g2-in-1);
      \draw (g2-out-0) to[out=0,in=180] (g3-in-0);
      \draw (g3-out-0) to[out=0,in=180] (g4-in-0);
    \end{tikzpicture}
\end{center}

        \end{figure}
        \item A diagram $c \in \mathbf{GPA}$

        \begin{figure}[H]
            \centering
            \begin{center}
    \;\begin{tikzpicture}[axpic]\node{$c$};\end{tikzpicture}\;
    \;\begin{tikzpicture}[axpic]\node{$=$};\end{tikzpicture}\;
    \begin{tikzpicture}[axpic]
        \pic[val=A] (g1) at (1.000,0.500) {relation};
        \pic (g2) at (2.500,0.500) {ge};
        \draw (g1-out-0) to[out=0,in=180] (g2-in-0);
        \pic[val=B] (g3) at (4.000,0.500) {comatrix};
        \draw (g2-out-0) to[out=0,in=180] (g3-in-0);
    \end{tikzpicture} \\
    \vspace{3mm}

    \;\begin{tikzpicture}[axpic]\node{$\overset{(zeroeq)}{=}$};\end{tikzpicture}\;
    \begin{tikzpicture}[axpic]
      \pic (g1) at (0.500,0.800) {linear};
      \pic (z1) at (3.000,0.800) {cozero};
      \draw (g1-out-0) to[out=0,in=180] (z1-in-0);
      \coordinate (id_in_2) at (0,0.000);
      \coordinate (id_out_3) at (1.000,0.000);
      \draw (id_in_2) -- (id_out_3);
      \pic[val=A] (g2) at (2.000,0.000) {relation};
      \pic (g3) at (3.500,0.000) {ge};
      \pic[val=B] (g4) at (5.000,0.000) {comatrix};
      \draw (id_out_3) to[out=0,in=180] (g2-in-0);
      \draw (g2-out-0) to[out=0,in=180] (g3-in-0);
      \draw (g3-out-0) to[out=0,in=180] (g4-in-0);
    \end{tikzpicture}\\
    \vspace{3mm}

    \;\begin{tikzpicture}[axpic]\node{$\overset{(inf)}{=}$};\end{tikzpicture}\;
    \begin{tikzpicture}[axpic]
      \pic (g1) at (0.500,0.800) {linear};
      \pic (z0) at (1.750,0.800) {ge};
      \pic (z1) at (3.000,0.800) {cozero};
      \draw (g1-out-0) to[out=0,in=180] (z0-in-0);
      \draw (z0-out-0) to[out=0,in=180] (z1-in-0);
      \coordinate (id_in_2) at (0,0.000);
      \coordinate (id_out_3) at (1.000,0.000);
      \draw (id_in_2) -- (id_out_3);
      \pic[val=A] (g2) at (2.000,0.000) {relation};
      \pic (g3) at (3.500,0.000) {ge};
      \pic[val=B] (g4) at (5.000,0.000) {comatrix};
      \draw (id_out_3) to[out=0,in=180] (g2-in-0);
      \draw (g2-out-0) to[out=0,in=180] (g3-in-0);
      \draw (g3-out-0) to[out=0,in=180] (g4-in-0);
    \end{tikzpicture}\\
    \vspace{3mm}

    \;\begin{tikzpicture}[axpic]\node{$\overset{(matrix\; def.)}{=}$};\end{tikzpicture}\;
    \begin{tikzpicture}[axpic]
      \pic (g1) at (0.500,0.300) {linear};
      \coordinate (id_in_2) at (0,-0.300);
      \coordinate (id_out_3) at (1.000,-0.300);
      \draw (id_in_2) -- (id_out_3);
      \pic[val=A'] (g2) at (2.000,0.000) {2drelation};
      \pic (g3) at (3.500,0.000) {ge};
      \pic[val=B'] (g4) at (5.000,0.000) {comatrix};
      \draw (g1-out-0) to[out=0,in=180] (g2-in-0);
      \draw (id_out_3) to[out=0,in=180] (g2-in-1);
      \draw (g2-out-0) to[out=0,in=180] (g3-in-0);
      \draw (g3-out-0) to[out=0,in=180] (g4-in-0);
    \end{tikzpicture}
\end{center}

        \end{figure}
    \end{itemize}
    For the inductive step, we verify that for any two diagrams
    $c, d \in \mathbf{GLP}$ satisfying the inductive hypothesis,
    both horizontal and vertical composition preserve the normal form.
    Closure under vertical composition follows by stacking the two normal forms and
    collapsing the two independent $A_i, \mathrm{ge}, B_i$ blocks into a single
    block-diagonal matrix pair, as shown below:
    \begin{figure}[H]
        \centering
        \begin{center}
    \begin{tikzpicture}[axpic]
      \begin{scope}[shift={(0,1.700)}]
      \pic (g1) at (0.500,0.800) {linear};
      \coordinate (id_in_1) at (0,0.200);
      \coordinate (id_out_1) at (1.000,0.200);
      \draw (id_in_1) -- (id_out_1);
      \pic[val=A_1] (g2) at (2.000,0.500) {2drelation};
      \pic (g3) at (3.500,0.500) {ge};
      \pic[val=B_1] (g4) at (5.000,0.500) {comatrix};
      \draw (g1-out-0) to[out=0,in=180] (g2-in-0);
      \draw (id_out_1) to[out=0,in=180] (g2-in-1);
      \draw (g2-out-0) to[out=0,in=180] (g3-in-0);
      \draw (g3-out-0) to[out=0,in=180] (g4-in-0);
      \end{scope}

      \pic (h1) at (0.500,0.800) {linear};
      \coordinate (id_in_2) at (0,0.200);
      \coordinate (id_out_2) at (1.000,0.200);
      \draw (id_in_2) -- (id_out_2);
      \pic[val=A_2] (h2) at (2.000,0.500) {2drelation};
      \pic (h3) at (3.500,0.500) {ge};
      \pic[val=B_2] (h4) at (5.000,0.500) {comatrix};
      \draw (h1-out-0) to[out=0,in=180] (h2-in-0);
      \draw (id_out_2) to[out=0,in=180] (h2-in-1);
      \draw (h2-out-0) to[out=0,in=180] (h3-in-0);
      \draw (h3-out-0) to[out=0,in=180] (h4-in-0);
    \end{tikzpicture} \\
    \vspace{3mm}

    \;\begin{tikzpicture}[axpic]\node{$=$};\end{tikzpicture}\;
    \begin{tikzpicture}[axpic]
      \begin{scope}[shift={(0,1.700)}]
      \pic (g1) at (0.500,0.800) {linear};
      \coordinate (id_in_1) at (0,0.200);
      \coordinate (id_out_1) at (1.000,0.200);
      \draw (id_in_1) -- (id_out_1);
      \pic[val=A_1] (g2) at (2.000,0.500) {2drelation};
      \pic (g3) at (3.500,0.500) {ge};
      \pic[val=B_1] (g4) at (5.000,0.500) {comatrix};
      \draw (g1-out-0) to[out=0,in=180] (g2-in-0);
      \draw (id_out_1) to[out=0,in=180] (g2-in-1);
      \draw (g2-out-0) to[out=0,in=180] (g3-in-0);
      \draw (g3-out-0) to[out=0,in=180] (g4-in-0);
      \end{scope}

      \pic (h1) at (0.500,0.800) {linear};
      \coordinate (id_in_2) at (0,0.200);
      \coordinate (id_out_2) at (1.000,0.200);
      \draw (id_in_2) -- (id_out_2);
      \pic[val=A_2] (h2) at (2.000,0.500) {2drelation};
      \pic (h3) at (3.500,0.500) {ge};
      \pic[val=B_2] (h4) at (5.000,0.500) {comatrix};
      \draw (h1-out-0) to[out=0,in=180] (h2-in-0);
      \draw (id_out_2) to[out=0,in=180] (h2-in-1);
      \draw (h2-out-0) to[out=0,in=180] (h3-in-0);
      \draw (h3-out-0) to[out=0,in=180] (h4-in-0);

      \draw[dashed] (1.500,0.000) rectangle (2.500,2.700);
      \draw[dashed] (4.400,0.000) rectangle (5.600,2.700);
    \end{tikzpicture} \\
    \vspace{3mm}

    \;\begin{tikzpicture}[axpic]\node{$=$};\end{tikzpicture}\;
    \begin{tikzpicture}[axpic]
      \pic (g1) at (0.500,0.800) {linear};
      \coordinate (id_in_3) at (0,0.200);
      \coordinate (id_out_3) at (1.000,0.200);
      \draw (id_in_3) -- (id_out_3);
      \pic[val=A'] (g2) at (2.000,0.500) {2drelation};
      \pic (g3) at (3.500,0.500) {ge};
      \pic[val=B'] (g4) at (5.000,0.500) {comatrix};
      \draw (g1-out-0) to[out=0,in=180] (g2-in-0);
      \draw (id_out_3) to[out=0,in=180] (g2-in-1);
      \draw (g2-out-0) to[out=0,in=180] (g3-in-0);
      \draw (g3-out-0) to[out=0,in=180] (g4-in-0);
    \end{tikzpicture}
\end{center}

    \end{figure}
    For horizontal composition, the two normal forms are joined along the wire connecting
    $B_1$ to $A_2$; applying Fourier-Motzkin Elimination to this pair recovers a fresh
    matrix-$\mathrm{ge}$-comatrix triple, after which the two resulting runs of matrices,
    one in each direction, compose into a single matrix and a single comatrix, recovering
    the normal form:
    \begin{figure}[H]
        \centering
        \begin{center}
    % ---------------------------------------------------------------
    % Layout convention: 1-unit pics (linear, ge, mult, ...) have ports
    % at +-0.5, 2-unit pics (matrix, comatrix, relation, ...) at +-1.0.
    % Consecutive pics are placed so that out-port <= next in-port.
    % ---------------------------------------------------------------
    \begin{tikzpicture}[axpic]
        \pic (g1) at (0.500,0.800) {linear};
        \coordinate (id_in_1) at (0,0.200);
        \coordinate (id_out_1) at (1.000,0.200);
        \draw (id_in_1) -- (id_out_1);
        \pic[val=A_1] (g2) at (2.000,0.500) {2drelation};
        \pic (g3) at (3.500,0.500) {ge};
        \pic[val=B_1] (g4) at (5.000,0.500) {comatrix};
        \draw (g1-out-0) to[out=0,in=180] (g2-in-0);
        \draw (id_out_1) to[out=0,in=180] (g2-in-1);
        \draw (g2-out-0) to[out=0,in=180] (g3-in-0);
        \draw (g3-out-0) to[out=0,in=180] (g4-in-0);

        \pic (h1) at (6.500,0.800) {linear};
        \pic[val=A_2] (h2) at (8.000,0.500) {2drelation};
        \pic (h3) at (9.500,0.500) {ge};
        \pic[val=B_2] (h4) at (11.000,0.500) {comatrix};
        \draw (h1-out-0) to[out=0,in=180] (h2-in-0);
        \draw (g4-out-0) to[out=0,in=180] (h2-in-1);
        \draw (h2-out-0) to[out=0,in=180] (h3-in-0);
        \draw (h3-out-0) to[out=0,in=180] (h4-in-0);
    \end{tikzpicture} \\
    \vspace{3mm}

    \;\begin{tikzpicture}[axpic]\node{$=$};\end{tikzpicture}\;
    \begin{tikzpicture}[axpic]
        \pic (g1) at (0.500,0.800) {linear};
        \coordinate (id_in_1) at (0,0.200);
        \coordinate (id_out_1) at (1.000,0.200);
        \draw (id_in_1) -- (id_out_1);
        \pic[val=A_1] (g2) at (2.000,0.500) {2drelation};
        \pic (g3) at (3.500,0.500) {ge};
        \pic[val=B_1] (g4) at (5.000,0.500) {comatrix};
        \draw (g1-out-0) to[out=0,in=180] (g2-in-0);
        \draw (id_out_1) to[out=0,in=180] (g2-in-1);
        \draw (g2-out-0) to[out=0,in=180] (g3-in-0);
        \draw (g3-out-0) to[out=0,in=180] (g4-in-0);

        \pic (h1) at (6.500,1.200) {linear};
        \pic[val=A_{2,1}] (n1) at (8.000,1.200) {matrix};
        \pic[val=A_{2,2}] (n2) at (8.000,-0.200) {matrix};
        \pic (mult1) at (10.000,0.500) {multiplication};
        \pic (aff1) at (8.500,-1.300) {affine};
        \pic[val=s] (sc1) at (9.500,-1.300) {scalar};
        \pic (mult2) at (11.500,-0.200) {multiplication};
        \pic (h3) at (12.500,-0.200) {ge};
        \pic[val=B_2] (h4) at (14.000,-0.200) {comatrix};
        \draw (h1-out-0) to[out=0,in=180] (n1-in-0);
        \draw (g4-out-0) to[out=0,in=180] (n2-in-0);
        \draw (n1-out-0) to[out=0,in=180] (mult1-in-0);
        \draw (n2-out-0) to[out=0,in=180] (mult1-in-1);
        \draw (aff1-out-0) to[out=0,in=180] (sc1-in-0);
        \draw (mult1-out-0) to[out=0,in=180] (mult2-in-0);
        \draw (sc1-out-0) to[out=0,in=180] (mult2-in-1);
        \draw (mult2-out-0) to[out=0,in=180] (h3-in-0);
        \draw (h3-out-0) to[out=0,in=180] (h4-in-0);
    \end{tikzpicture} \\
    \vspace{3mm}

    \;\begin{tikzpicture}[axpic]\node{$=$};\end{tikzpicture}\;
    \begin{tikzpicture}[axpic]
        \pic (g1) at (0.500,0.800) {linear};
        \coordinate (id_in_1) at (0,0.200);
        \coordinate (id_out_1) at (1.000,0.200);
        \draw (id_in_1) -- (id_out_1);
        \pic[val=A_1] (g2) at (2.000,0.500) {2drelation};
        \pic (g3) at (3.500,0.500) {ge};
        \pic[val=B_1] (g4) at (5.000,0.500) {comatrix};
        \draw (g1-out-0) to[out=0,in=180] (g2-in-0);
        \draw (id_out_1) to[out=0,in=180] (g2-in-1);
        \draw (g2-out-0) to[out=0,in=180] (g3-in-0);
        \draw (g3-out-0) to[out=0,in=180] (g4-in-0);

        \pic (h1) at (6.500,1.200) {linear};
        \pic[val=A_{2,1}] (n1) at (8.000,1.200) {matrix};
        \pic[val=A_{2,2}] (n2) at (8.000,-0.200) {matrix};
        \pic (h3) at (9.500,-0.200) {ge};
        \pic (mult1) at (11.000,0.500) {multiplication};
        \pic (aff1) at (9.500,-1.300) {affine};
        \pic[val=s] (sc1) at (10.500,-1.300) {scalar};
        \pic (mult2) at (12.500,-0.200) {multiplication};
        \pic[val=B_2] (h4) at (14.000,-0.200) {comatrix};
        \draw (h1-out-0) to[out=0,in=180] (n1-in-0);
        \draw (g4-out-0) to[out=0,in=180] (n2-in-0);
        \draw (n2-out-0) to[out=0,in=180] (h3-in-0);
        \draw (n1-out-0) to[out=0,in=180] (mult1-in-0);
        \draw (h3-out-0) to[out=0,in=180] (mult1-in-1);
        \draw (aff1-out-0) to[out=0,in=180] (sc1-in-0);
        \draw (mult1-out-0) to[out=0,in=180] (mult2-in-0);
        \draw (sc1-out-0) to[out=0,in=180] (mult2-in-1);
        \draw (mult2-out-0) to[out=0,in=180] (h4-in-0);
    \end{tikzpicture} \\
    \vspace{3mm}

    \;\begin{tikzpicture}[axpic]\node{$\overset{(FME)}{=}$};\end{tikzpicture}\;
    \begin{tikzpicture}[axpic]
        % first normal form, centred on y = 1.5
        \pic (g1) at (0.500,1.800) {linear};
        \coordinate (id_in_2) at (0,1.200);
        \coordinate (id_out_3) at (1.000,1.200);
        \draw (id_in_2) -- (id_out_3);
        \pic[val=A_1] (g4) at (2.000,1.500) {2drelation};
        \draw (g1-out-0) to[out=0,in=180] (g4-in-0);
        \draw (id_out_3) to[out=0,in=180] (g4-in-1);
        \pic[val=C] (g5) at (4.000,1.500) {matrix};
        \draw (g4-out-0) to[out=0,in=180] (g5-in-0);
        % middle row
        \pic (g8) at (5.500,1.500) {ge};
        \pic[val=D] (g9) at (7.000,1.500) {comatrix};
        \draw (g5-out-0) to[out=0,in=180] (g8-in-0);
        \draw (g8-out-0) to[out=0,in=180] (g9-in-0);
        % top row
        \pic (g6) at (5.500,2.700) {linear};
        \pic[val=A_{2,1}] (g7) at (7.000,2.700) {matrix};
        \draw (g6-out-0) to[out=0,in=180] (g7-in-0);
        \pic (g10) at (9.000,2.100) {multiplication};
        \draw (g7-out-0) to[out=0,in=180] (g10-in-0);
        \draw (g9-out-0) to[out=0,in=180] (g10-in-1);
        % bottom row
        \pic (g11) at (6.500,0.400) {affine};
        \pic[val=s] (g12) at (7.500,0.400) {scalar};
        \draw (g11-out-0) to[out=0,in=180] (g12-in-0);
        % merge
        \pic (g14) at (10.500,1.250) {multiplication};
        \draw (g10-out-0) to[out=0,in=180] (g14-in-0);
        \draw (g12-out-0) to[out=0,in=180] (g14-in-1);
        \pic[val=B_2] (g15) at (12.000,1.250) {comatrix};
        \draw (g14-out-0) to[out=0,in=180] (g15-in-0);
    \end{tikzpicture} \\
    \vspace{3mm}

    \;\begin{tikzpicture}[axpic]\node{$=$};\end{tikzpicture}\;
    \begin{tikzpicture}[axpic]
        \pic (g1) at (0.500,1.800) {linear};
        \coordinate (id_in_2) at (0,1.200);
        \coordinate (id_out_3) at (1.000,1.200);
        \draw (id_in_2) -- (id_out_3);
        \pic[val=A_1] (g4) at (2.000,1.500) {2drelation};
        \draw (g1-out-0) to[out=0,in=180] (g4-in-0);
        \draw (id_out_3) to[out=0,in=180] (g4-in-1);
        \pic[val=C] (g5) at (4.000,1.500) {matrix};
        \draw (g4-out-0) to[out=0,in=180] (g5-in-0);
        % middle row
        \pic (g9) at (5.500,1.500) {ge};
        \draw (g5-out-0) to[out=0,in=180] (g9-in-0);
        % top row
        \pic (g6) at (2.500,2.700) {linear};
        \pic[val=A_{2,1}] (g7) at (4.000,2.700) {matrix};
        \pic[val=D] (g8) at (6.000,2.700) {matrix};
        \draw (g6-out-0) to[out=0,in=180] (g7-in-0);
        \draw (g7-out-0) to[out=0,in=180] (g8-in-0);
        \pic (g12) at (8.000,2.100) {multiplication};
        \draw (g8-out-0) to[out=0,in=180] (g12-in-0);
        \draw (g9-out-0) to[out=0,in=180] (g12-in-1);
        \pic[val=D] (g13) at (9.500,2.100) {comatrix};
        \draw (g12-out-0) to[out=0,in=180] (g13-in-0);
        % bottom row
        \pic (g14) at (7.000,0.400) {affine};
        \pic[val=s] (g15) at (8.000,0.400) {scalar};
        \draw (g14-out-0) to[out=0,in=180] (g15-in-0);
        % merge
        \pic (g17) at (11.500,1.250) {multiplication};
        \draw (g13-out-0) to[out=0,in=180] (g17-in-0);
        \draw (g15-out-0) to[out=0,in=180] (g17-in-1);
        \pic[val=B_2] (g18) at (13.000,1.250) {comatrix};
        \draw (g17-out-0) to[out=0,in=180] (g18-in-0);
    \end{tikzpicture} \\
    \vspace{3mm}

    \;\begin{tikzpicture}[axpic]\node{$\overset{(\text{antipode})}{=}$};\end{tikzpicture}\;
    \begin{tikzpicture}[axpic]
        \pic (g1) at (0.500,1.800) {linear};
        \coordinate (id_in_2) at (0,1.200);
        \coordinate (id_out_3) at (1.000,1.200);
        \draw (id_in_2) -- (id_out_3);
        \pic[val=A_1] (g4) at (2.000,1.500) {2drelation};
        \draw (g1-out-0) to[out=0,in=180] (g4-in-0);
        \draw (id_out_3) to[out=0,in=180] (g4-in-1);
        \pic[val=C] (g5) at (4.000,1.500) {matrix};
        \draw (g4-out-0) to[out=0,in=180] (g5-in-0);
        % middle row
        \pic (g9) at (5.500,1.500) {ge};
        \draw (g5-out-0) to[out=0,in=180] (g9-in-0);
        % top row
        \pic (g6) at (2.500,2.700) {linear};
        \pic[val=A_{2,1}] (g7) at (4.000,2.700) {matrix};
        \pic[val=D] (g8) at (6.000,2.700) {matrix};
        \draw (g6-out-0) to[out=0,in=180] (g7-in-0);
        \draw (g7-out-0) to[out=0,in=180] (g8-in-0);
        \pic (g12) at (8.000,2.100) {multiplication};
        \draw (g8-out-0) to[out=0,in=180] (g12-in-0);
        \draw (g9-out-0) to[out=0,in=180] (g12-in-1);
        \pic[val=D] (g13) at (9.500,2.100) {comatrix};
        \draw (g12-out-0) to[out=0,in=180] (g13-in-0);
        % split
        \pic (g14) at (11.000,2.100) {comultiplication};
        \draw (g13-out-0) to[out=0,in=180] (g14-in-0);
        \pic[val=B_2] (g15) at (12.500,2.600) {comatrix};
        \draw (g14-out-0) to[out=0,in=180] (g15-in-0);
        \draw (g15-out-0) -- ++(0.250,0);
        \pic (g16) at (12.000,1.600) {antipode};
        \pic[val=s] (g17) at (13.000,1.600) {coscalar};
        \pic (g18) at (13.750,1.600) {coaffine};
        \draw (g14-out-1) to[out=0,in=180] (g16-in-0);
        \draw (g16-out-0) to[out=0,in=180] (g17-in-0);
    \end{tikzpicture} \\
    \vspace{3mm}

    \;\begin{tikzpicture}[axpic]\node{$=$};\end{tikzpicture}\;
    \begin{tikzpicture}[axpic]
        \pic (g1) at (0.500,1.800) {linear};
        \coordinate (id_in_2) at (0,1.200);
        \coordinate (id_out_3) at (1.000,1.200);
        \draw (id_in_2) -- (id_out_3);
        \pic[val=A_1] (g4) at (2.000,1.500) {2drelation};
        \draw (g1-out-0) to[out=0,in=180] (g4-in-0);
        \draw (id_out_3) to[out=0,in=180] (g4-in-1);
        \pic[val=C] (g5) at (4.000,1.500) {matrix};
        \draw (g4-out-0) to[out=0,in=180] (g5-in-0);
        % top row
        \pic (g6) at (2.500,2.700) {linear};
        \pic[val=A_{2,1}] (g7) at (4.000,2.700) {matrix};
        \pic[val=D] (g8) at (6.000,2.700) {matrix};
        \draw (g6-out-0) to[out=0,in=180] (g7-in-0);
        \draw (g7-out-0) to[out=0,in=180] (g8-in-0);
        \pic (g13) at (8.000,2.100) {multiplication};
        \draw (g8-out-0) to[out=0,in=180] (g13-in-0);
        \draw (g5-out-0) to[out=0,in=180] (g13-in-1);
        \pic (g14) at (9.000,2.100) {ge};
        \draw (g13-out-0) to[out=0,in=180] (g14-in-0);
        \pic[val=D] (g15) at (10.500,2.100) {comatrix};
        \draw (g14-out-0) to[out=0,in=180] (g15-in-0);
        % split
        \pic (g16) at (12.000,2.100) {comultiplication};
        \draw (g15-out-0) to[out=0,in=180] (g16-in-0);
        \pic[val=B_2] (g17) at (13.500,2.600) {comatrix};
        \draw (g16-out-0) to[out=0,in=180] (g17-in-0);
        \draw (g17-out-0) -- ++(0.250,0);
        \pic (g18) at (13.000,1.600) {antipode};
        \pic[val=s] (g19) at (14.000,1.600) {coscalar};
        \pic (g20) at (14.750,1.600) {coaffine};
        \draw (g16-out-1) to[out=0,in=180] (g18-in-0);
        \draw (g18-out-0) to[out=0,in=180] (g19-in-0);
    \end{tikzpicture} \\
    \vspace{3mm}

    \;\begin{tikzpicture}[axpic]\node{$\overset{(\text{matrix comp.})}{=}$};\end{tikzpicture}\;
    \begin{tikzpicture}[axpic]
        \pic (g1) at (0.500,1.300) {linear};
        \coordinate (id_in_2) at (0,0.700);
        \coordinate (id_out_3) at (1.000,0.700);
        \draw (id_in_2) -- (id_out_3);
        \pic[val=A''] (g4) at (2.000,1.000) {2drelation};
        \draw (g1-out-0) to[out=0,in=180] (g4-in-0);
        \draw (id_out_3) to[out=0,in=180] (g4-in-1);
        \pic (g5) at (3.500,1.000) {ge};
        \draw (g4-out-0) to[out=0,in=180] (g5-in-0);
        \pic[val=B''] (g6) at (5.000,1.000) {corelation};
        \draw (g5-out-0) to[out=0,in=180] (g6-in-0);
    \end{tikzpicture}

\end{center}

    \end{figure}
\end{proof}

The following corollary reveals the geometric meaning of the normal form:
the $\mathbf{GPA}$ part of every $\mathbf{GLP}$ diagram encodes the
epigraph of the corresponding value function. This is the key structural
insight connecting the syntax to the semantic equivalence of
Proposition~\ref{prop:equalepi}.

\begin{corollary}[Epigraph Form]\label{cor:epigraph form}
    Given a diagram $d \in \mathbf{GLP}$ in normal form, it can always
    be rewritten as a diagram $d'$ of the form:
    \begin{figure}[H]
        \centering
        \begin{tikzpicture}
  \pic (g1) at (0.500,1.000) {linear};
  \pic (g2) at (1.750,1.000) {comultiplication};
  \draw (g1-out-0) to[out=0,in=180] (g2-in-0);
  \pic[val=A] (g3) at (3.750,1.000) {2drelation};
  \draw (g2-out-0) to[out=0,in=180] (g3-in-0);
  \draw (g2-out-1) to[out=0,in=180] (g3-in-1);
  \pic (g4) at (5.250,1.000) {ge};
  \draw (g3-out-0) to[out=0,in=180] (g4-in-0);
  \pic (g5) at (6.250,1.000) {comultiplication};
  \draw (g4-out-0) to[out=0,in=180] (g5-in-0);
  \pic[val=B] (g6) at (7.750,1.500) {comatrix};
  \draw (g5-out-0) to[out=0,in=180] (g6-in-0);
  \draw (g6-out-0) -- ++(1.000,0);
  \pic (g7) at (7.250,0.500) {antipode};
  \draw (g5-out-1) to[out=0,in=180] (g7-in-0);
  \pic[val=C] (g8) at (8.750,0.500) {comatrix};
  \draw (g7-out-0) to[out=0,in=180] (g8-in-0);
  \draw[dotted] (1.150,-0.050) rectangle (9.500,2.050);
\end{tikzpicture}

    \end{figure}
    where the semantics of the dotted diagram is
    \[\left\{(t,y)\mid \exists\, x,\; A'x + a \ge [B\; -C]\, y,\; t\ge x,\; y =
    \begin{bmatrix} y_1 \\ y_2 \end{bmatrix}\right\}\].
    This is the epigraph of the original diagram $d$
    \begin{align*}
        [d](y_1, y_2) &= \inf_{z} \; z \quad \text{s.t. } A z + a \ge B y_1 - C y_2 \\
        \mathrm{epi}([d]) &= \left\{(t, y_1, y_2) \mid \exists\, z \le t,\;
            A z + a \ge B y_1 - C y_2 \right\}
    \end{align*}
    such that $A = [A' \;C]$. We call this the \emph{epigraph form} of a linear program.
\end{corollary}

\begin{proof}
    Starting from $d$ in normal form, split the relation generator along
    $A = [A' \; C]$ and pull the resulting $C$-block through the $\inlinegen{ge}$
    generator using the antipode axiom:
    \begin{center}
\begin{tikzpicture}[axpic]
  \pic (g1) at (0.500,1.300) {linear};
  \coordinate (id_in_2) at (0,0.700);
  \coordinate (id_out_3) at (1.000,0.700);
  \draw (id_in_2) -- (id_out_3);
  \pic[val=A] (g4) at (2.000,1.000) {2drelation};
  \draw (g1-out-0) to[out=0,in=180] (g4-in-0);
  \draw (id_out_3) to[out=0,in=180] (g4-in-1);
  \pic (g5) at (3.500,1.000) {ge};
  \draw (g4-out-0) to[out=0,in=180] (g5-in-0);
  \pic[val=B] (g6) at (5.000,1.000) {comatrix};
  \draw (g5-out-0) to[out=0,in=180] (g6-in-0);
\end{tikzpicture} \\
\vspace{3mm}

\;\begin{tikzpicture}[axpic]\node{$=$};\end{tikzpicture}\;
\begin{tikzpicture}[axpic]
  \pic (g1) at (0.500,1.500) {linear};
  \coordinate (id_in_2) at (0,0.500);
  \coordinate (id_out_3) at (1.000,0.500);
  \draw (id_in_2) -- (id_out_3);
  \pic[val=A] (g4) at (2.000,1.500) {relation};
  \pic[val=C] (g5) at (2.000,0.500) {matrix};
  \draw (g1-out-0) to[out=0,in=180] (g4-in-0);
  \draw (id_out_3) to[out=0,in=180] (g5-in-0);
  \pic (g6) at (3.500,1.000) {multiplication};
  \draw (g4-out-0) to[out=0,in=180] (g6-in-0);
  \draw (g5-out-0) to[out=0,in=180] (g6-in-1);
  \pic (g7) at (4.500,1.000) {ge};
  \draw (g6-out-0) to[out=0,in=180] (g7-in-0);
  \pic[val=B] (g8) at (6.000,1.000) {comatrix};
  \draw (g7-out-0) to[out=0,in=180] (g8-in-0);
\end{tikzpicture} \\
\vspace{3mm}

\;\begin{tikzpicture}[axpic]\node{$=$};\end{tikzpicture}\;
\begin{tikzpicture}[axpic]
  \pic (g1) at (0.500,1.000) {linear};
  \pic[val=A] (g2) at (2.000,1.000) {relation};
  \draw (g1-out-0) to[out=0,in=180] (g2-in-0);
  \pic (g3) at (3.500,1.000) {ge};
  \draw (g2-out-0) to[out=0,in=180] (g3-in-0);
  \pic (g4) at (4.500,1.000) {comultiplication};
  \draw (g3-out-0) to[out=0,in=180] (g4-in-0);
  \pic[val=B] (g5) at (6.000,1.500) {comatrix};
  \draw (g4-out-0) to[out=0,in=180] (g5-in-0);
  \draw (g5-out-0) -- ++(1.000,0);
  \pic (g6) at (5.500,0.500) {antipode};
  \draw (g4-out-1) to[out=0,in=180] (g6-in-0);
  \pic[val=C] (g7) at (7.000,0.500) {comatrix};
  \draw (g6-out-0) to[out=0,in=180] (g7-in-0);
\end{tikzpicture} \\
\vspace{3mm}

\;\begin{tikzpicture}[axpic]\node{$=$};\end{tikzpicture}\;
\begin{tikzpicture}[axpic]
  \pic (g1) at (0.500,1.000) {linear};
  \pic (g2) at (1.750,1.000) {comultiplication};
  \draw (g1-out-0) to[out=0,in=180] (g2-in-0);
  \pic[val=A] (g3) at (3.750,1.000) {2drelation};
  \draw (g2-out-0) to[out=0,in=180] (g3-in-0);
  \draw (g2-out-1) to[out=0,in=180] (g3-in-1);
  \pic (g4) at (5.250,1.000) {ge};
  \draw (g3-out-0) to[out=0,in=180] (g4-in-0);
  \pic (g5) at (6.250,1.000) {comultiplication};
  \draw (g4-out-0) to[out=0,in=180] (g5-in-0);
  \pic[val=B] (g6) at (7.750,1.500) {comatrix};
  \draw (g5-out-0) to[out=0,in=180] (g6-in-0);
  \draw (g6-out-0) -- ++(1.000,0);
  \pic (g7) at (7.250,0.500) {antipode};
  \draw (g5-out-1) to[out=0,in=180] (g7-in-0);
  \pic[val=C] (g8) at (8.750,0.500) {comatrix};
  \draw (g7-out-0) to[out=0,in=180] (g8-in-0);
  \draw[dotted] (1.150,-0.050) rectangle (9.500,2.050);
\end{tikzpicture}
\end{center}

    The resulting diagram is precisely the claimed form $d'$, with the
    dotted box exhibiting $\mathrm{epi}([d])$ as a $\mathbf{GPA}$ diagram.
\end{proof}

\section{Presentation results for \texorpdfstring{$\mathbf{LPRel}$}{LPRel}}

With the syntax and its normal form in place, we are now ready to prove
the three presentation properties for $\mathbf{GLP}$.

\begin{lemma}[Expressivity]
    For every $f : m \to n$ in $\mathbf{LPRel}$ there exists a
    $\Sigma$-term $s : m \to n$ in
    $\mathrm{Hom}_{\mathrm{FreeP}_{\mathbf{GLP}}}(m,n)$
    such that $[s] = f$.
\end{lemma}

\begin{proof}
    Since every polyhedral function is the maximum of a finite family of
    affine functions, it suffices to show that every piecewise affine
    convex function is expressible in $\mathbf{GLP}$. For an arbitrary
    such function $f(y) = \max_i R_i(y)$, where $\{R_i\}_{i=1,\dots,n}$
    is a family of affine functions, the following diagram provides the
    required syntactic representative:
    \begin{figure}[H]
        \centering
        \begin{tikzpicture}[axpic]
    % top row: objective and constraint matrix
    \pic (g1) at (0.500,1.500) {linear};
    \pic[val=c^\top] (g2) at (1.500,1.500) {coscalar};
    \draw (g1-out-0) to[out=0,in=180] (g2-in-0);
    \pic[val=A] (g3) at (3.000,1.500) {matrix};
    \draw (g2-out-0) to[out=0,in=180] (g3-in-0);
    % bottom row: affine offset
    \pic (g4) at (2.500,0.500) {affine};
    \coordinate (a5) at (4.000,0.500);
    \draw (g4-out-0) -- (a5);
    % merge, constrain, output
    \pic (g6) at (4.500,1.000) {multiplication};
    \draw (g3-out-0) to[out=0,in=180] (g6-in-0);
    \draw (a5)       to[out=0,in=180] (g6-in-1);
    \pic (g7) at (5.500,1.000) {ge};
    \draw (g6-out-0) to[out=0,in=180] (g7-in-0);
    \pic[val=B] (g8) at (7.000,1.000) {comatrix};
    \draw (g7-out-0) to[out=0,in=180] (g8-in-0);
    \node[anchor=west] at ($(g8-out-0)+(0.15,0)$)
        {$
            \llbracket = \rrbracket \; y \mapsto \inf_x c^\top x
            \ \text{s.t.}\
            Ax + b \ge By$};
\end{tikzpicture}

    \end{figure}
\end{proof}

\begin{lemma}[Soundness]
   If $s = t$ in $\mathrm{FreeP}_{\mathbf{GLP}}$, then $[s] = [t]$
   in $\mathbf{LPRel}$.
\end{lemma}

\begin{proof}
    All axioms inherited from $\mathbf{GPA}$ are already sound by the
    results of the previous section. It suffices to verify that the new
    axioms involving the LP generator preserve the semantics, which
    follows by direct computation:
    \begin{figure}[H]
        \centering
        \begin{tikzpicture}
    \node at (-1.5,0) {\textit{(linear.ge)}};
    \pic at (0,0) {linear};
    \node at (1.2,0) {$\llbracket = \rrbracket$};
    \node at (2.6,0) {$(\bullet, y) \mapsto y$};
    \node at (3.6,0) {$ = $};
    \node at (5.2,0) {$(\bullet, y) \mapsto \inf_{x \ge y} x$};
    \node at (7.4,0) {$\llbracket = \rrbracket$};
    \pic at (8.3,0) {linear};
    \draw (8.8,0) to (9.3,0);
    \pic at (9.5,0) {ge};
\end{tikzpicture}
\begin{tikzpicture}
    \node at (-1.5,0) {(\textit{lin.sum})};
    \pic (g1) at (0.000,0.25) {linear};
    \pic (g2) at (0.000,-0.25) {linear};
    \node at (1.2,0) {$\llbracket = \rrbracket$};
    \node at (3.5,0) {$(\bullet,\left(\begin{smallmatrix} y_1 \\ y_2 \end{smallmatrix}\right)) \mapsto y_1 + y_2$};
    \node at (5.8,0) {$\llbracket = \rrbracket$};
    \pic at (6.6,0) {linear};
    \draw (7.1,0) to (7.6,0);
    \pic at (7.8,0) {comultiplication};
\end{tikzpicture}
\vspace{6mm}
\begin{tikzpicture}
    \node at (-1.5,0) {(\textit{lin.zero})};
    \pic at (0.0,0) {linear};
    \draw (.2,0) to (0.5,0);
    \pic at (1.0,0) {cozero};
    \node at (1.7,0) {$\llbracket = \rrbracket$};
    \node at (3.9,0) {$(\bullet, \bullet) \mapsto \inf_{x = 0} x = 0$};
    \node at (6.3,0) {$\llbracket = \rrbracket$};
    \node[draw, dotted] at (7.2,0) {};
\end{tikzpicture}

    \end{figure}
\end{proof}

The final result of the thesis is the completeness theorem for
$\mathbf{GLP}$. The proof brings together all the main tools developed
throughout: the epigraph form of Corollary~\ref{cor:epigraph form}, the
shared epigraph property of Proposition~\ref{prop:equalepi}, the
completeness of $\mathbf{GPA}$, and the FME-based normal form.

\begin{lemma}[Completeness]
   If $[c] = [d]$ in $\mathbf{LPRel}$, then $c = d$ in
   $\mathrm{FreeP}_{\mathbf{GLP}}$.
\end{lemma}

\begin{proof}
    Let $v_1 = [c]$ and $v_2 = [d]$. By assumption $v_1 = v_2$, so
    Proposition~\ref{prop:equalepi} gives
    $\mathrm{epi}(v_1) = \mathrm{epi}(v_2) =: E$.

    By Corollary~\ref{cor:epigraph form}, both $c$ and $d$ admit
    epigraph forms $c = \mushroom ; c^*$ and $d = \mushroom ; d^*$
    with $[c^*] = [d^*] = E$ in $\mathbf{PolyRel}$. Since $c^*$ and $d^*$
    are $\mathbf{GPA}$ diagrams with the same denotation, completeness of
    $\mathbf{GPA}$ gives $c^* = d^*$ in $\mathrm{FreeP}_{\mathbf{GPA}}$;
    in particular both reduce to the same $\mathbf{GPA}$ normal form.
    Composing this common diagram with the leftover $\mushroom$ collapses
    $c$ and $d$ to the same $\mathbf{GLP}$ diagram:
    \begin{figure}[H]
        \centering
        \begin{center}
    \begin{tikzpicture}[axpic]
      \node at (0.5,1.3) {$c$};
      \pic (g1) at (0.500,0.800) {linear};
      \coordinate (id_in_1) at (0,0.200);
      \coordinate (id_out_1) at (1.000,0.200);
      \draw (id_in_1) -- (id_out_1);
      \pic[val=A_1] (g2) at (2.000,0.500) {2drelation};
      \pic (g3) at (3.500,0.500) {ge};
      \pic[val=B_1] (g4) at (5.000,0.500) {corelation};
      \draw (g1-out-0) to[out=0,in=180] (g2-in-0);
      \draw (id_out_1) to[out=0,in=180] (g2-in-1);
      \draw (g2-out-0) to[out=0,in=180] (g3-in-0);
      \draw (g3-out-0) to[out=0,in=180] (g4-in-0);

      \begin{scope}[shift={(7.000,0.0)}]
      \node at (0.5,1.3) {$d$};
      \pic (h1) at (0.500,0.800) {linear};
      \coordinate (id_in_2) at (0,0.200);
      \coordinate (id_out_2) at (1.000,0.200);
      \draw (id_in_2) -- (id_out_2);
      \pic[val=A_2] (h2) at (2.000,0.500) {2drelation};
      \pic (h3) at (3.500,0.500) {ge};
      \pic[val=B_2] (h4) at (5.000,0.500) {corelation};
      \draw (h1-out-0) to[out=0,in=180] (h2-in-0);
      \draw (id_out_2) to[out=0,in=180] (h2-in-1);
      \draw (h2-out-0) to[out=0,in=180] (h3-in-0);
      \draw (h3-out-0) to[out=0,in=180] (h4-in-0);
      \end{scope}
    \end{tikzpicture} \\
    \vspace{3mm}
    {\footnotesize (epigraph form, Cor.~\ref{cor:epigraph form})} \\
    \vspace{1mm}

    \begin{tikzpicture}[axpic]
      \pic (g1) at (0.500,0.500) {linear};
      \pic (g2) at (1.500,0.500) {ge};
      \pic (g3) at (2.500,0.500) {comultiplication};
      \pic[val=A_1] (g4) at (4.500,0.500) {2dmatrix};
      \pic (g5) at (6.000,0.500) {ge};
      \pic[val=B_1] (g6) at (7.500,0.500) {corelation};
      \draw (g1-out-0) to[out=0,in=180] (g2-in-0);
      \draw (g2-out-0) to[out=0,in=180] (g3-in-0);
      \draw (g3-out-0) to[out=0,in=180] (g4-in-0);
      \draw (g3-out-1) to[out=0,in=180] (g4-in-1);
      \draw (g4-out-0) to[out=0,in=180] (g5-in-0);
      \draw (g5-out-0) to[out=0,in=180] (g6-in-0);
      \draw[dotted] (1.100,-0.150) rectangle (8.300,1.150);
      \node at (4.700,-0.400) {\scriptsize $c^*$};
    \end{tikzpicture} \\
    \vspace{1mm}

    \begin{tikzpicture}[axpic]
      \pic (h1) at (0.500,0.500) {linear};
      \pic (h2) at (1.500,0.500) {ge};
      \pic (h3) at (2.500,0.500) {comultiplication};
      \pic[val=A_2] (h4) at (4.500,0.500) {2dmatrix};
      \pic (h5) at (6.000,0.500) {ge};
      \pic[val=B_2] (h6) at (7.500,0.500) {corelation};
      \draw (h1-out-0) to[out=0,in=180] (h2-in-0);
      \draw (h2-out-0) to[out=0,in=180] (h3-in-0);
      \draw (h3-out-0) to[out=0,in=180] (h4-in-0);
      \draw (h3-out-1) to[out=0,in=180] (h4-in-1);
      \draw (h4-out-0) to[out=0,in=180] (h5-in-0);
      \draw (h5-out-0) to[out=0,in=180] (h6-in-0);
      \draw[dotted] (1.100,-0.150) rectangle (8.300,1.150);
      \node at (4.700,-0.400) {\scriptsize $d^*$};
    \end{tikzpicture} \\
    \vspace{3mm}
    {\footnotesize (the dotted boxes $c^*, d^*$ are $\mathbf{GPA}$ diagrams with $[c^*]=[d^*]=E$;
      by $\mathbf{GPA}$ completeness both reduce to the same normal form)} \\
    \vspace{1mm}

    \begin{tikzpicture}[axpic]
      \pic (g1) at (0.500,0.500) {linear};
      \pic[val=A^*] (g2) at (2.000,0.500) {relation};
      \pic (g3) at (3.500,0.500) {ge};
      \pic[val=B^*] (g4) at (5.000,0.500) {comatrix};
      \draw (g1-out-0) to[out=0,in=180] (g2-in-0);
      \draw (g2-out-0) to[out=0,in=180] (g3-in-0);
      \draw (g3-out-0) to[out=0,in=180] (g4-in-0);
      \draw[dotted] (1.100,0.000) rectangle (5.750,1.000);

      \node at (6.550,0.500) {$=$};

      \begin{scope}[shift={(7.200,0.0)}]
      \pic (h1) at (0.500,0.500) {linear};
      \pic[val=A^*] (h2) at (2.000,0.500) {relation};
      \pic (h3) at (3.500,0.500) {ge};
      \pic[val=B^*] (h4) at (5.000,0.500) {comatrix};
      \draw (h1-out-0) to[out=0,in=180] (h2-in-0);
      \draw (h2-out-0) to[out=0,in=180] (h3-in-0);
      \draw (h3-out-0) to[out=0,in=180] (h4-in-0);
      \draw[dotted] (1.100,0.000) rectangle (5.750,1.000);
      \end{scope}
    \end{tikzpicture} \\
    \vspace{3mm}
    {\footnotesize (composing the common $\mathbf{GPA}$ normal form with the leftover
      \inlinegen{linear} yields the same $\mathbf{GLP}$ diagram)} \\
    \vspace{1mm}

    \begin{tikzpicture}[axpic]
      \pic (g1) at (0.500,0.500) {linear};
      \pic[val=A^*] (g2) at (2.000,0.500) {relation};
      \pic (g3) at (3.500,0.500) {ge};
      \pic[val=B^*] (g4) at (5.000,0.500) {comatrix};
      \draw (g1-out-0) to[out=0,in=180] (g2-in-0);
      \draw (g2-out-0) to[out=0,in=180] (g3-in-0);
      \draw (g3-out-0) to[out=0,in=180] (g4-in-0);
    \end{tikzpicture}
\end{center}

    \end{figure}
    Therefore $c = \mushroom ; c^* = \mushroom ; d^* = d$.
\end{proof}

As a consequence of the three lemmas above:

\begin{theorem}
    $\mathbf{LPRel}$ is presented by $\mathrm{FreeP}_{\mathbf{GLP}}$.
\end{theorem}
\begin{proof}
    By the preceding three lemmas.
\end{proof}

Finally, we've built, through the addition of one generator and three axioms,
an elegant and complete diagrammatical theory for Parametric Linear Programming.
We argue that, such result, not just unlocks a new way to interact with LP problems in a algebraic
and computational way, but also provide a new perspective to optimization through the lens of
a compositional framework. Whereas optimization, namely LP problems, are lifted above
the settings of functional thinking and achieves a new status as a relational-theoretic
field embeded into the Extended-Lawvere Quantale category.

The \textbf{GLP} theory, in turn, has a natural extension to an even more expressive theory
when enriched with the gray arrow generator from \textbf{GQA}. We believe that the class
of quadratic programs could be, just as was done in this work, fully axiomatized by such
broader SMIT. However, such studies are left for future works, along with other possible applications
of current results.
